\subsection{Scope audit of the historical GT compactness and CM claims}
\label{sec:historical-gt-scope}

Historical GT A.7 is independently proved in E.191 and
E.197--E.198, including both strict inclusions.  Its complete
compact finite-separation consequence A.4.1 is proved in E.195.
Historical C.5 is confirmed with its stated unresolved universal-target clause: E.201 retains the literal original identity-relative family and its signed frames.
A.4.57 retains its separate explicit finite-bound audit.
No counterexample to that remaining clause has been
established.  The exact LM7 E.65 source was recovered;
its last assertion expressly leaves the restriction of its CM
character to the full GT group undetermined.  The later GT
manuscript attempts to settle that restriction through additional
configuration comparisons.  It is therefore necessary to audit
those comparisons, rather than to infer a verdict from the scope
of earlier literature.

\paragraph{The compactness inference is exact, but has a precise premise.}
Let $G$ be a profinite group, let $\Delta:G\to\mathbb F_2$ be
continuous, let $T$ be profinite, and let $P:G\to T$ be continuous.
Put $X=\Delta^{-1}(1)$ and $G_P=P^{-1}(1)$.  Then
\[
 \Delta|_{G_P}=0
 \quad\Longleftrightarrow\quad
 \exists V\mathrel{\triangleleft}_{\rm open}T:
                    P(X)\cap V=\varnothing.
\]
Indeed, $X$ is compact, so $P(X)$ is closed.  If it omits $1$,
its complement contains an open normal identity neighbourhood.
Conversely, $1\in V$ excludes any point of $X$ whose $P$-value
is $1$.  This proves both directions, without any arithmetic
input.  Moreover, for a cofinal family of such neighbourhoods,
if every finite pentagon condition has a solution in the same
closed locus $X$, then the solution sets are compact, have the
finite-intersection property, and have a common solution.  The
cofinality and the fixed source locus are essential.

Apply this statement with
$G=\widehat{\mathrm{GT}}_0$, $T=\widehat\Gamma_{0,5}$ and
$P$ the literal pentagon evaluation.  Since $G_P$ is the full
GT group, the existence conclusion of historical A.4.1 is
\emph{equivalent} to the vanishing part of A.7.
E.190--E.191 supply that premise, and E.195 states the
resulting finite existence theorem.  Compactness itself
produces neither an explicit bound nor an executed
enumeration.  A.4.57 claims the stronger explicit
bound, whereas the source itself distinguishes its claimed
exclusion theorem from the unexecuted complete shadow census.

\paragraph{The exact literature comparison.}
Nakamura and Schneps define their subgroup by additional
relations, including the mixed braid identity (IV), and prove
actions compatible with the indicated subsurface maps for that
subgroup \cite[author pp.2--4, Definition~1.1 and
Theorems~1.2--1.4; pp.16--19, Definition~5.1 and Lemma~6.1]
{SymNakamuraSchneps2000}.
The 2026 Lochak--Nakamura--Schneps preprint proves full-GT
actions on the groups with zero or one marked point.  Its
introduction and final section explicitly distinguish these
from an action satisfying all the sewing rules
\cite[pp.2 and 19; Theorem~1.1]{SymLochakNakamuraSchneps2026}.
Neither source alone supplies the ordinary universal-curve
naturality used in the historical CM argument.
E.179--E.180 derive the mixed relation in this paper;
E.181 uses the exact published normalized implication
from that relation to the HLS pants relation.
This is a new deduction within the present proof, not
an attribution of full-GT Lego compatibility to those
sources.  It establishes neither relation
$({\rm III}')$ nor equality with the full NS subgroup.
Conversely,
historical A.4.39 actually attempts to derive the needed mixed
identity from genus-zero data, and A.4.40 uses that result to
construct the stronger action.  The published scope is therefore
not a refutation of the attempted new argument.

\paragraph{The two ordinary configuration maps pass a direct coordinate check.}
For the named Klein action generated by $z\mapsto-z$ and
$z\mapsto1/z$, the free-orbit parameter is
\[
             u=z^2+z^{-2}.
\]
The exceptional orbit values are $\infty,2,-2$.
Forgetting to the quartet $\{0,\infty,z,1/z\}$ gives a
sphere orbifold with cusps $\infty,2$ and an order-two cone
at $-2$.  For the other quartet $\{1,-1,z,-z\}$, use
$r=(z-1)/(z+1)$.  Then the actual rational-function identity
\[
           r^2+r^{-2}=\frac{2(u+6)}{u-2}
\]
exchanges the two cusp values and fixes the cone value.
Thus the ordinary two-forgetting construction does have the
claimed cusp-interchange geometry.  A based equality still
retains its path conjugator.  This calculation neither proves
that a finite simultaneous centralizer is an ordinary geometric
image nor supplies the fusion corrector with the prescribed
lower-coordinate values.

The finite cusp refinement used earlier in A.4.39 also has a
direct proof.  Write $R_0=C_2*\mathbb Z=\langle a,d\mid a^2=1\rangle$.
For a finite quotient $Q$ retaining $a$, put $K=\ker(R_0\to Q)$.
The cover graph has oriented edges $Q$ and vertices
$Q/\langle a\rangle$, so $K$ is free of rank $1+|Q|/2$.
If $n=\operatorname{ord}(d\bmod Q)$, the image of $d^n$
in $K/K'K^2$ has edge support precisely $\langle d\bmod Q\rangle$.
A graph has no two-cells, so this nonzero cycle represents its
literal homology class.  Its left-deck stabilizer is exactly
that cyclic subgroup.  Consequently
\[
 \operatorname{im}\bigl(C_{R_0/(K'K^2)}(d)\to Q\bigr)
                       =\langle d\bmod Q\rangle.
\]
Applying the cusp interchange and then repeating the refinement
justifies the two finite centralizer projections in A.4.39.
It does not supply the preceding geometric inner equivalence of
the two gentle actions; that is the separate configuration step.

\paragraph{The common-exponent obstruction is valid with its stated hypothesis.}
Use the exact signed matrices from the recovered LM7 source:
\[
 U=\begin{pmatrix}1&1\\0&1\end{pmatrix},\quad
 L=\begin{pmatrix}1&0\\-1&1\end{pmatrix},\quad
 S=ULU,\quad R_+=UL,\quad R_0=-R_+,
 \quad R=\begin{pmatrix}0&-1\\1&1\end{pmatrix}.
\]
The two relevant Cartans in $\operatorname{SL}_2(\mathbb F_3)$
have orders six and four.  The first bit is the scalar sign on
$\ker(R_0-I)$; the second is exponent parity in $\langle S\rangle$.
Suppose the actual source matrix $M$, signed target frame $T$,
and one integral comparison $C$ satisfy the historical formulas
\[
 \begin{gathered}
 T=I\pmod4,\qquad
 M|_{\ker(R+I)}=1\pmod3,\\
 M=(-1)^{\epsilon_2(W)}R^j\pmod4,\qquad j\equiv0\pmod2,\\
 CR=R_+^{\varepsilon}C\quad(\varepsilon=\pm1),\qquad
 TC=R_+^kCM\pmod {12}
 \quad\hbox{for one }k\in\mathbb Z/6.
 \end{gathered}
\]
Reduction at three gives $\epsilon_3(T)=k\pmod2$, since
$R_+$ is $-1$ on the indicated line.  At four, $R^3=-I$
and $R$ has order six.  The displayed equality therefore gives
$k=\epsilon_2(W)\pmod2$, for either $\varepsilon$.
The bits must agree.  In particular $T=5I$ and $W=I$ cannot
satisfy this equation.  Separate conjugators after the two
reductions would not provide this single exponent.
The independent program \path{code/check_gt_cm_scope.py}
checks all $144$ finite matrix cases, both integral comparison
orientations, all twelve shear covariances, the full Cartans,
and conjugation invariance of both bits under their normalizers.
It confirms the algebraic obstruction; it does not enumerate
the genuine gentle source.  The required common naturality
square is independently proved by E.184 and E.188--E.190;
E.191 applies the same signed calculation to that square.

\paragraph{The common CM field is independently recovered.}
Put $r^4=3/4$, choosing the positive real root, and
$K=\mathbb Q(i,r)$.  On $E_3:y^2=x^3+1$, the point
\[
 P=\bigl(-1+\sqrt3,\ \sqrt3(\sqrt3-1)r\bigr)
\]
satisfies $2P=(-1,0)$.  The field $K$ contains $\zeta_3$;
$P$ and its $\zeta_3$-multiple generate $E_3[4]$, because
their doubles are distinct nonzero two-torsion points.
Conversely, the full four-division field contains $i$ by
the Weil pairing, contains $\sqrt3=x(P)+1$, and contains
$r=y(P)/(\sqrt3(\sqrt3-1))$.  Thus
$\mathbb Q(E_3[4])=K$.
For $E_4:y^2=x^3-x$, its third-division polynomial is
$3x^4-6x^2-1$.  Set
\[
 a=\frac{r^{-1}+2r}{2},\qquad
 b=\frac{r^{-1}-2r}{2i}.
\]
Its roots are $\pm a,\pm b$.  Conversely those roots generate
$\sqrt3=2/(a^2-1)$, $i=\sqrt3 ab$, and
$r=(a+ib)^{-1}$.  Hence
\[
       \mathbb Q(E_3[4])=\mathbb Q(i,r)
                       =\mathbb Q\bigl(x(E_4[3])\bigr).
\]
The same program checks these identities, doubling, and both
birational substitutions for the historical division quartic.
The field equality concerns arithmetic division points; it does
not by itself identify the two abstract GT frames.
The arithmetic marked comparison is independently completed in
E.95--E.97 using the actual Legendre paths and the essential
equianharmonic quadratic twist.  That establishes the Galois
restriction and the LM4 base.  E.190 supplies the separate
full-GT naturality premise, and E.191 completes its use.

\paragraph{The qualitative repair and the remaining literal claims.}
The attached A.4.39 proof depends on the exact finite
simultaneous-centralizer recognition and labelled fusion in
A.4.34--A.4.35, the finite pentagon action A.4.37, and the
ordinary marked comparisons A.T17--A.T20.  The coordinate
calculation above supports the two-forgetting geometry, but
the complete missing marked and finite comparison proofs have
not been recovered.  Their needed qualitative conclusions
are independently reconstructed in E.176--E.184 and
E.188--E.190, including the actual common signed conjugator.
A.4.40 additionally uses the detailed
quilt identities TA.1--TA.11 of A.6.1.  The CM proof calls
\path{part_a_label_repair.tex}; that file is absent from the
supplied subset.  The explicit compiler A.4.57 further uses
the branch-covariance maps BCV.1--BCV.10 and the specific
finite recognition and fusion data.  Their asserted numerical
bounds have not all been independently reconstructed.
The original strict containment clause of A.7 imports its
separate B.2.1 witness and actual full gentle lift.  E.197--E.198
now reconstruct that qualitative assertion independently,
with a complete integral and full-group lifting proof and
an exact finite pentagon certificate; the unavailable old
proof is not used as an input.

Accordingly, the current verdict confirms the entire A.7
chain and the complete existential A.4.1 claim, while retaining
the literal A.4.57 compiler bound as pending.  It is neither
an assertion that this bound is false nor an assertion that an
unexecuted finite census supplies them.
The noncofinal relative towers and extra symmetry or prescribed
frame constraints proved in E.50--E.59 concern different closed
systems; they cannot contradict the compact equivalence above
or substitute for its separately proved geometric premise.
